% 英文长摘要的中文对照稿。论点、数字与文献与
% qingdao-posts-extended-abstract.tex 一致，不另增内容。
% 不使用 LNCS：llncs 不排中文。本文件只供对照阅读。
\documentclass[11pt]{article}
\usepackage{fontspec}
\usepackage{geometry}
\geometry{a4paper, top=2.5cm, bottom=2.5cm, left=2.7cm, right=2.7cm}
\setmainfont{PingFang SC}[
  BoldFont={PingFang SC Semibold}
]
\newfontfamily\enit{texgyretermes}[
  Path={/Users/a123/Library/TinyTeX/texmf-dist/fonts/opentype/public/tex-gyre/},
  Extension=.otf,
  UprightFont=texgyretermes-regular,
  BoldFont=texgyretermes-bold,
  ItalicFont=texgyretermes-italic,
  BoldItalicFont=texgyretermes-bolditalic
]
\XeTeXlinebreaklocale "zh"
\XeTeXlinebreakskip=0pt plus 0.5pt
\linespread{1.35}\selectfont
\pagestyle{plain}
\renewcommand{\abstractname}{摘要}
\renewcommand{\refname}{参考文献}

\begin{document}

\begin{center}
{\Large\bfseries 谁在表征目的地，他们评价什么？\\[0.4em]
关于青岛的社交媒体帖子\par}
\vspace{1.1em}
{\large 匿名作者\par}
\vspace{0.7em}
{\normalsize 双盲评审，作者姓名与单位从略。\par}
\vspace{0.8em}
{\small 英文长摘要的中文对照稿。论点、数字与文献与英文稿一致。\par}
\end{center}

\vspace{0.6em}
\noindent\textbf{摘要。}
目的地营销者和研究者常常把社交媒体帖子直接当作游客所持目的地形象的替身。本文将帖子理解为被生产出来的话语，而不是静止的认知形象，考察机构账号和个人创作者如何呈现同一个目的地青岛。帖子清洗了 19,099 篇。放弃 B 站低质量帖子后，四个平台供正式研究使用的为 10,178 篇。对 100 篇（476 个片段）的试点分析中，实用的行程规划信息占多数（占片段的 58\%），情感与感官表达较少（30\%）。对核心吸引物的评价绝大多数为正向（166 条中的 155 条）；价格、服务、可达性和拥挤等运营方面则正负大致相当（正向 20 条，负向 19 条）。下一步考察机构创作者和个人创作者是否会从同一批资源投射出不同的目的地形象。

\vspace{0.8em}
\noindent\textbf{关键词：}目的地形象；用户生成内容；内容创作者；内容分析。

\section{背景}
目的地形象是一个人对某地持有的累积信念和印象，不能还原成内容频次的简单合计（Tasci, Gartner, \& Cavusgil, 2007）。社交媒体帖子为形象形成提供材料（Gartner, 1994），也可以编码认知属性和情感属性（Garay, 2019），但网上的一篇帖子并不等于访客内心持有的认知形象。既有研究对照官方营销和游客生成内容时，主要用词频看主题是否出现（Sun, Tang, \& Liu, 2021）。目的地在话语上怎样被表征，因此仍知道得不多：帖子承担什么实用功能，评价哪些属性，创作者身份又怎样塑造叙述。

\section{目的}
本研究考察机构账号和个人创作者如何表征青岛，看他们使用什么内容策略、评价哪些资源和属性，以及这些要素如何在离散的文本片段里共现。我们的假设是：即便写到的是同一批目的地资源，机构帖子也更稳定地偏向不加批评的推广性表征；个人创作者则在总体正向的叙述里加入实用性的保留、批评或建议。

\section{方法}
方法是理论导向的结构化内容分析（Krippendorff, 2018）。片段通过对帖子正文做语义切分得到。每个片段编码内容策略、所指向的目的地资源、被评价属性，以及评价的方向和强度。账号元数据与文本编码分开记录。创作者分组标准在查看组间比较结果之前已经预先登记。清洗模型在 1,300 条标注样本上训练后，得到供正式研究使用的 10,178 篇帖子，并因放弃 B 站低质量帖子而改为四个平台。试点样本由两名经过培训的人类编码员独立编码，评价方向的编码员间信度为\mbox{Cohen's {\enit κ} = 0.82}。分歧通过协商讨论解决，并据此确定最终标签。

\section{现阶段结果}
\noindent\textbf{筛选与目的地相关的内容。}
19,099 篇帖子经过清洗。放弃 B 站低质量帖子后，四个平台上供正式研究使用的为 10,178 篇。按关键词匹配的原始帖子计数，会高估实际使用的帖子。

\vspace{0.7em}
\noindent\textbf{帖子在做什么。}
试点样本是随机抽取的 100 篇帖子，共 476 个片段。实用信息（行程、价格、做法建议）出现在 58\% 的片段中，感官或情感呈现占 30\%，背景知识占 11\%，互动号召占 3\%。大约一半帖子给出明确推荐。

\vspace{0.7em}
\noindent\textbf{帖子在评价什么。}
两名人类编码员在 216 条评价上，对目标属性和评价方向都一致。核心吸引物质量、活动和目的地氛围共 166 条（正向 155，负向 10）。运营因素——价格、服务、设施、可达性、拥挤和管理——共 44 条，正负大致相当（正向 20，负向 19）。

\vspace{0.7em}
\noindent\textbf{类别在何处分不清。}
分歧最常出在两处：到访一处地标算活动还是景观，以及某类建筑算景观还是服务设施。修订后的编码簿因此把景观性建筑吸引物和功能性接待设施明确分开。

\section{意义}
对目的地营销组织来说，这次试点把总情感分数盖住的运营差别拆开了。正向评价主要落在目的地固有的资源禀赋上，抱怨则集中在可以管理的运营接触点。能够据以采取管理行动的信号，大约只占全部评价的五分之一。帖子也主要是实用的出行规划工具，而不是单纯的情感性个人抒发。如果随后的全样本分析确认，个人创作者会吸收并投射运营层面的批评，机构账号则把这些批评略去，那么两组人是在用同一批资源构建出不同的目的地。这与会议主题 \mbox{Hype and Return} 一致。本方法表明，计算自动化在何处带来效率，严格的人工辨别在何处仍然不可替代。

\section*{参考文献}
{\raggedright
\begin{list}{}{\leftmargin=1.6em \itemindent=-1.6em \itemsep=0.45em \parsep=0pt}
\item Garay, L. (2019). \#Visitspain. Breaking down affective and cognitive attributes in the social media construction of the tourist destination image. {\enit\itshape Tourism Management Perspectives, 32}, 100560.
\item Gartner, W. C. (1994). Image formation process. {\enit\itshape Journal of Travel \& Tourism Marketing, 2}(2--3), 191--216.
\item Krippendorff, K. (2018). {\enit\itshape Content analysis: An introduction to its methodology} (4th ed.). Sage.
\item Sun, W., Tang, S., \& Liu, F. (2021). Examining perceived and projected destination image: A social media content analysis. {\enit\itshape Sustainability, 13}(6), 3354.
\item Tasci, A. D. A., Gartner, W. C., \& Cavusgil, S. T. (2007). Conceptualization and operationalization of destination image. {\enit\itshape Journal of Hospitality \& Tourism Research, 31}(2), 194--223.
\end{list}
}

\end{document}
